\section{Matlis Principal Parts and the Continuous Cotangent Module}
\label{sec:app-matlis-principal-parts}

This appendix centralizes the reduced Matlis calculation used in
Theorem~\ref{thm:lane-principal-parts}.  Its stable core is the
following: the \(m\)-adic continuous dual of the reduced Hamiltonian
Lie algebra is the reduced scalar-annihilator inside the top
local-cohomology module; the residue pairing gives the coadjoint
action; and no ordinary polynomial descendant module realizes this
coadjoint representation.  The residue pairing is the universal dual
module normalization for this lane.  The Fourier--Rees construction below
then gives the strongest flat positive bridge that exists: a genuine
\(\mathcal D_\hbar\)-module bridge after Fourier conjugation of the
polarization, together with an exact obstruction theorem for any bridge
that tries to preserve the original Hamiltonian action.  Over
\(\C[[\hbar]]\) its target is the Fourier delta Rees lattice inside local
cohomology, not the full standard Cech lattice; the full principal-part
lattice appears only after tensoring with \(\C((\hbar))\).  It does not
identify \(\Psi_{a,b}\) with \(\rho_{a,b}\) as \(\mathfrak h\)-modules.
The source conventions are
Matlis duality
\cite{matlis-injective}, Grothendieck--Hartshorne local cohomology
\cite{hartshorne-local-cohomology}, and the power-series residue model
of Kunz--Cox--Dickenstein, Chapter~5
\cite{kunz-residues-duality}.

\begin{defn}[Topology and reduced Hamiltonians]
\label{def:app-matlis-topology}
  Put
  \[
    R=\C[[z_1,z_2]],\qquad \mathfrak m=(z_1,z_2),\qquad
    \mathfrak h=R/\C\cdot1 .
  \]
  The topology on \(R\) is the \(\mathfrak m\)-adic topology.  The
  topology on \(\mathfrak h\) is the quotient topology, equivalently
  the topology transported through the zero-constant section
  \[
    \mathfrak h\simeq \mathfrak m,\qquad
    \overline f\longmapsto f-f(0).
  \]
  A fundamental system of neighbourhoods of zero is
  \[
    F^N\mathfrak h=\mathfrak m^N,\qquad N\geq1,
  \]
  via this splitting.  The Lie bracket is the Poisson bracket
  \[
    \{f,g\}=\partial_{z_1}f\,\partial_{z_2}g
      -\partial_{z_2}f\,\partial_{z_1}g
  \]
  followed by projection modulo constants.
\end{defn}

\begin{lem}[Continuous dual]
\label{lem:app-matlis-continuous-dual}
  The continuous dual of \(\mathfrak h\), with \(\C\) discrete, is
  \[
    \mathfrak h^\vee_{\mathrm{cont}}
      =\bigcup_{N\geq1}\bigl(\mathfrak h/F^N\mathfrak h\bigr)^\vee
      =\bigoplus_{\substack{a,b\geq0\\a+b>0}}\C\,\delta_{a,b},
  \]
  where
  \[
    \delta_{a,b}(z_1^m z_2^n)=\delta_{a,m}\delta_{b,n}.
  \]
\end{lem}

\begin{proof}
  A linear functional on the \(\mathfrak m\)-adic vector space
  \(\mathfrak h\simeq\mathfrak m\) is continuous exactly when its
  kernel contains \(F^N\mathfrak h=\mathfrak m^N\) for some \(N\).
  Such a functional factors through the finite-dimensional quotient
  \(\mathfrak m/\mathfrak m^N\), hence is a finite linear combination
  of Taylor coefficient functionals.  Conversely every finite linear
  combination of the displayed \(\delta_{a,b}\) kills all sufficiently
  high powers of \(\mathfrak m\), hence is continuous.
\end{proof}

\begin{defn}[Reduced Matlis principal parts]
\label{def:app-matlis-principal-parts}
  The top local cohomology module of \(R\) is written in Cech
  Laurent form as
  \[
    H^2_{\mathfrak m}(R)\,dz_1dz_2
      =\bigoplus_{a,b\geq0}
        \C\,z_1^{-a-1}z_2^{-b-1}\,dz_1dz_2 .
  \]
  Equivalently this is the Cech quotient
  \[
    R[z_1^{-1},z_2^{-1}]/(R[z_1^{-1}]+R[z_2^{-1}])
  \]
  after multiplication by \(dz_1dz_2\); the displayed direct sum is
  the standard representative basis in which both exponents are
  strictly negative.
  The reduced principal-part module is the \(\C\)-linear residue annihilator
  of constants,
  \[
    \mc P=\operatorname{Ann}_{\mathrm{res}}(\C\cdot1)
      =\bigoplus_{\substack{a,b\geq0\\a+b>0}}\C\,\rho_{a,b},
      \qquad
    \rho_{a,b}=z_1^{-a-1}z_2^{-b-1}\,dz_1dz_2 .
  \]
  The residue pairing is
  \[
    \langle \rho_{a,b},z_1^m z_2^n\rangle
      =\operatorname{Res}_{z_1=z_2=0}
        z_1^{m-a-1}z_2^{n-b-1}\,dz_1dz_2
      =\delta_{a,m}\delta_{b,n}.
  \]
\end{defn}

\begin{prop}[Matlis realization of the continuous dual]
\label{prop:app-matlis-realization}
  The residue pairing induces a continuous vector-space isomorphism
  \[
    \mc P\xrightarrow{\sim}\mathfrak h^\vee_{\mathrm{cont}},
    \qquad
    \rho_{a,b}\longmapsto\delta_{a,b}.
  \]
  It is the scalar-reduced restriction of the Matlis evaluation
  pairing for the complete regular local ring \(R\): after the
  trivialization \(\omega_R=R\,dz_1dz_2\),
  \(H^2_{\mathfrak m}(R)\,dz_1dz_2\cong E_R(\C)\).  The subspace
  \(\mc P=\operatorname{Ann}_{\mathrm{res}}(\C\cdot1)\) is the
  \(\C\)-linear scalar-annihilator dual to
  \(\mathfrak h=R/\C\cdot1\).  The \(R\)-module structure on
  \(E_R(\C)\) does not restrict to \(\mc P\); multiplication can hit the
  excluded scalar residue line.  The Hamiltonian coadjoint action below
  is transported separately by the residue pairing and is the structure
  used by the brane cotangent sector.
\end{prop}

\begin{proof}
  The displayed residue formula identifies each basis vector
  \(\rho_{a,b}\), \(a+b>0\), with the continuous Taylor coefficient
  functional \(\delta_{a,b}\).  Lemma~\ref{lem:app-matlis-continuous-dual}
  shows that these coefficient functionals form the whole continuous
  dual.  The exclusion of \(\rho_{0,0}\) is exactly the condition that
  the functional vanish on the constant Hamiltonian \(1\), which has
  been removed from \(\mathfrak h\).
\end{proof}

\begin{prop}[Residue coadjoint action]
\label{prop:app-matlis-coadjoint-action}
  Under the isomorphism of
  Proposition~\ref{prop:app-matlis-realization}, the coadjoint action
  of \(f\in\mathfrak h\) is the principal-part residue action
  \[
    f\cdot\rho=\pi_{\mathrm{pp}}\{f,\rho\},
  \]
  where \(\pi_{\mathrm{pp}}\) projects Laurent two-forms to the span
  of the principal parts.  On monomials,
  \[
    z_1^p z_2^q\cdot\rho_{a,b}
      =-(pb-qa+p-q)\rho_{a-p+1,b-q+1},
  \]
  with the convention that any term with a negative index is zero.
\end{prop}

\begin{proof}
  Hamiltonian vector fields preserve the volume form \(dz_1dz_2\).
  Hence the residue pairing satisfies integration by parts,
  \[
    \langle \{f,\rho\},g\rangle
      =-\langle \rho,\{f,g\}\rangle ,
  \]
  which is exactly the coadjoint sign convention
  \((\operatorname{ad}^{\vee}_f\lambda)(g)=-\lambda(\{f,g\})\).

  For \(f=z_1^p z_2^q\) and
  \(\rho_{a,b}=z_1^{-a-1}z_2^{-b-1}\,dz_1dz_2\),
  \[
  \begin{aligned}
    \{f,z_1^{-a-1}z_2^{-b-1}\}
      &=
      pz_1^{p-1}z_2^q\,(-(b+1))z_1^{-a-1}z_2^{-b-2}  \\
      &\quad
      -qz_1^p z_2^{q-1}\,(-(a+1))z_1^{-a-2}z_2^{-b-1} \\
      &=-(pb-qa+p-q)z_1^{p-a-2}z_2^{q-b-2}.
  \end{aligned}
  \]
  Multiplication by \(dz_1dz_2\) gives
  \(\rho_{a-p+1,b-q+1}\) precisely when the two new indices are
  nonnegative; otherwise the principal-part projection gives zero.
\end{proof}

\begin{thm}[Residue pairing rigidity]
\label{thm:app-matlis-residue-rigidity}
  Up to multiplication by one scalar in \(\C^\times\), the residue
  pairing is the unique nondegenerate continuous pairing
  \[
    \mc P\times\mathfrak h\longrightarrow\C
  \]
  of total degree zero which is equivariant for the coadjoint action
  above.  Fixing the scalar by
  \[
    \langle\rho_{a,b},z_1^m z_2^n\rangle
      =\delta_{a,m}\delta_{b,n}
  \]
  is therefore the universal dual-module normalization for the
  principal-part cotangent sector.
\end{thm}

\begin{proof}
  Let \(\beta\) be another such pairing and write
  \(\beta_{a,b;m,n}=\beta(\rho_{a,b},z_1^m z_2^n)\).  Total-degree
  zero gives \(\beta_{a,b;m,n}=0\) unless \(a+b=m+n\).
  Equivariance for \(z_1z_2\) gives the second constraint
  \(a-b=m-n\).  Together these force \((a,b)=(m,n)\), so
  \[
    \beta_{a,b;m,n}=c_{a,b}\delta_{a,m}\delta_{b,n}.
  \]
  Equivariance for \(z_1\) gives
  \(c_{a,b+1}=c_{a,b}\), because
  \(z_1\cdot\rho_{a,b}=-(b+1)\rho_{a,b+1}\) and
  \(\{z_1,z_1^m z_2^n\}=n z_1^m z_2^{n-1}\).  Equivariance for
  \(z_2\) gives \(c_{a+1,b}=c_{a,b}\).  All coefficients with
  \(a+b>0\) are therefore a single scalar \(c\), and
  nondegeneracy forces \(c\neq0\).
\end{proof}

\begin{thm}[Polynomial descendant obstruction for principal parts]
\label{thm:app-matlis-polynomial-descendant-obstruction}
  Let \(M\) be a polynomial boundary descendant module on which the
  linear Hamiltonians \(z_1,z_2\) act locally nilpotently.  Then
  \(\mc P\) is not isomorphic to any nonzero submodule of \(M\) as an
  \(\mathfrak h\)-module.  In particular the polynomial classes
  \(\Psi_{a,b}\) cannot be identified with the principal parts
  \(\rho_{a,b}\) by an ordinary \(\mathfrak h\)-module isomorphism.
\end{thm}

\begin{proof}
  In any ordinary polynomial boundary model, the linear Hamiltonians
  act as constant coefficient Hamiltonian vector fields; iterating
  either \(z_1\) or \(z_2\) eventually exhausts the polynomial degree.
  Thus these actions are locally nilpotent.

  On \(\mc P\), Proposition~\ref{prop:app-matlis-coadjoint-action}
  gives
  \[
    z_1\cdot\rho_{a,b}=-(b+1)\rho_{a,b+1},\qquad
    z_2\cdot\rho_{a,b}=(a+1)\rho_{a+1,b}.
  \]
  Hence
  \[
    z_1^N\cdot\rho_{a,b}
      =(-1)^N(b+1)(b+2)\cdots(b+N)\rho_{a,b+N}\neq0
  \]
  for every \(N\), and similarly for \(z_2\).  Local nilpotence is
  preserved by module isomorphism and by restriction to submodules.
  Therefore no polynomial descendant module with locally nilpotent
  linear Hamiltonians realizes \(\mc P\).
\end{proof}

\begin{thm}[Universal Fourier--Rees construction and exact obstruction]
\label{thm:app-matlis-rees-fourier-bridge}
  Let
  \[
    \mathcal D_\hbar
    =
    \C[[\hbar]]
      \langle z_1,z_2,\partial_1,\partial_2\rangle
      \big/
      \bigl([\partial_i,z_j]=\hbar\delta_{ij},
      [z_i,z_j]=[\partial_i,\partial_j]=0\bigr)
  \]
  be the Rees Weyl algebra.  Let
  \[
    \mathcal O_\hbar=\C[z_1,z_2][[\hbar]]
      \cong\mathcal D_\hbar/(\partial_1,\partial_2)
  \]
  be the polynomial Rees module, and let
  \[
    \Delta^{\mathsf F}_\hbar=
      \mathcal D_\hbar/(z_1,z_2)
  \]
  be the Fourier delta Rees module, generated by a vector \(e_0\)
  killed by \(z_1,z_2\).  The Cech local-cohomology lattice
  \[
    \Delta^{\mathrm{std}}_\hbar=
    H^2_{\mathfrak m}(\C[[z_1,z_2]])\,dz_1dz_2[[\hbar]]
  \]
  carries the standard action in which \(z_i\) multiplies and
  \(\partial_i=\hbar\partial_{z_i}\).  Sending
  \(e_0\mapsto\rho_{0,0}=z_1^{-1}z_2^{-1}dz_1dz_2\) realizes
  \(\Delta^{\mathsf F}_\hbar\) as the Rees sublattice
  \[
    \iota_{\mathsf F}(\Delta^{\mathsf F}_\hbar)
    =
    \bigoplus_{a,b\geq0}
      \C[[\hbar]]\,
      (-1)^{a+b}\hbar^{a+b}a!b!\,\rho_{a,b}
    \subset
    \Delta^{\mathrm{std}}_\hbar .
  \]
  This inclusion becomes an isomorphism after tensoring with
  \(\C((\hbar))\), but it is not an isomorphism of the two
  \(\C[[\hbar]]\)-lattices.  The Rees Fourier
  automorphism
  \[
    \mathsf F(z_i)=\partial_i,\qquad
    \mathsf F(\partial_i)=-z_i
  \]
  identifies the Fourier twist \(\mathsf F(\mathcal O_\hbar)\) with
  \(\Delta^{\mathsf F}_\hbar\).  This is the flat positive bridge: the
  action of \(P\in\mathcal D_\hbar\) on the polynomial source is carried
  to the action of \(\mathsf F(P)\) on the Fourier delta target.  On basis
  labels in the Cech
  realization,
  \[
    z_1^az_2^b
    \longmapsto
    \partial_1^a\partial_2^b e_0
    \overset{\iota_{\mathsf F}}{\longmapsto}
    \partial_1^a\partial_2^b\rho_{0,0}
    =
    (-1)^{a+b}\hbar^{a+b}a!b!\,\rho_{a,b}.
  \]
  After inverting \(\hbar\), this gives a filtered Fourier--Rees
  identification of label spaces between nonconstant polynomial labels and
  the reduced principal parts \(\rho_{a,b}\), \(a+b>0\), but only after the
  Fourier twist of action just stated.  Over the Rees lattice before
  inverting \(\hbar\), the construction records the degree-\(a+b\) shift
  by the explicit factor \(\hbar^{a+b}\).  The special fibre of the
  abstract Rees module keeps the labels
  \(\partial_1^a\partial_2^be_0\); the Cech inclusion into
  \(\Delta^{\mathrm{std}}_\hbar\) sends the positive-degree labels into
  \(\hbar\Delta^{\mathrm{std}}_\hbar\).  Thus the comparison is an
  associated-graded label match, not an isomorphism of ordinary
  local-cohomology lattices.

  The obstruction for the original Hamiltonian action is exact.  The
  Hamiltonian action on polynomial descendants is by vector fields
  \(X_f=\{f,-\}\) on \(\mathcal O_\hbar((\hbar))\).  In the Rees Weyl
  algebra the normalized linear operators are
  \(\hbar X_{z_1}=\partial_2\) and
  \(\hbar X_{z_2}=-\partial_1\).  Their Fourier transforms are
  \(\mathsf F(\hbar X_f)\), whereas the Matlis cotangent action on
  \(\mc P((\hbar))\) is still \(\hbar X_f\), acting on Laurent
  principal parts and projected to the reduced principal-part span.
  Already for the linear Hamiltonians,
  \[
    \hbar X_{z_1}=\partial_2,\qquad
    \hbar X_{z_2}=-\partial_1,
  \]
  but
  \[
    \mathsf F(\hbar X_{z_1})=-z_2,\qquad
    \mathsf F(\hbar X_{z_2})=z_1 .
  \]
  The transformed linear operators are locally nilpotent on
  \(\Delta^{\mathsf F}_\hbar\); the Rees-normalized coadjoint operators
  \[
    \hbar X_{z_1}\rho_{a,b}=-(b+1)\hbar\rho_{a,b+1},\qquad
    \hbar X_{z_2}\rho_{a,b}=(a+1)\hbar\rho_{a+1,b}
  \]
  are not locally nilpotent on any nonzero reduced principal part after
  inverting \(\hbar\).  Dividing by \(\hbar\) gives the same conclusion
  for the unnormalized Hamiltonian vector fields.  Hence the
  Fourier--Rees construction cannot identify the polynomial descendant
  module with \(\mc P\) as a module over the same reduced Hamiltonian Lie
  algebra.  Thus any universal bridge must choose one of two precise
  alternatives: Fourier-conjugate the action as above, or keep the original
  Matlis coadjoint action and accept the same-action obstruction.
\end{thm}

\begin{proof}
  The displayed formulas for \(\mathsf F\) preserve the Rees Weyl
  relations, since
  \[
    [\mathsf F(\partial_i),\mathsf F(z_j)]
      =[-z_i,\partial_j]=\hbar\delta_{ij}.
  \]
  The polynomial module \(\mathcal O_\hbar\) is cyclic with generator
  \(1\) annihilated by \(\partial_1,\partial_2\).  After Fourier twist
  the cyclic generator is annihilated by \(z_1,z_2\), so the twisted
  module is \(\mathcal D_\hbar/(z_1,z_2)=\Delta^{\mathsf F}_\hbar\).
  In the Cech model, \(z_i\rho_{0,0}=0\), and applying
  \(\partial_i=\hbar\partial_{z_i}\) to \(\rho_{0,0}\) generates the
  Rees lattice displayed in the theorem.  It generates the full
  standard local-cohomology lattice only after \(\hbar\) is inverted.
  Therefore
  \(\mathsf F(\mathcal O_\hbar)\cong\Delta^{\mathsf F}_\hbar\), with the
  Cech realization \(\iota_{\mathsf F}\).

  Since \(\partial_i\) acts as \(\hbar\partial_{z_i}\) on Laurent
  representatives,
  \[
    \partial_1^a\partial_2^b\rho_{0,0}
    =
    \hbar^{a+b}
    \partial_{z_1}^a\partial_{z_2}^b
      (z_1^{-1}z_2^{-1}dz_1dz_2)
    =
    (-1)^{a+b}\hbar^{a+b}a!b!\rho_{a,b}.
  \]
  Removing the scalar label \(a=b=0\) gives the reduced label comparison
  after the stated Rees normalization.  Since the factors \(a!b!\) and
  \((-1)^{a+b}\) are units over \(\C\), the only genuine Rees-lattice
  shift is the power \(\hbar^{a+b}\).

  For the obstruction, compare the two linear actions.  On polynomial
  descendants the Rees-normalized operators
  \(\hbar X_{z_1}=\partial_2\) and
  \(\hbar X_{z_2}=-\partial_1\) are locally nilpotent; after inverting
  \(\hbar\), the same is true for \(X_{z_1}\) and \(X_{z_2}\).  Fourier
  sends the Rees-normalized operators to multiplication by \(-z_2\) and
  \(z_1\) on \(\Delta^{\mathsf F}_\hbar\); multiplication by \(z_i\)
  lowers pole order and is again locally nilpotent.  The coadjoint
  Matlis action is not the Fourier-transformed action.  It keeps the
  same vector fields, now acting on Laurent principal parts:
  \[
    \hbar X_{z_1}\rho_{a,b}
      =\hbar\partial_{z_2}\rho_{a,b}
      =-(b+1)\hbar\rho_{a,b+1},\qquad
    \hbar X_{z_2}\rho_{a,b}
      =-\hbar\partial_{z_1}\rho_{a,b}
      =(a+1)\hbar\rho_{a+1,b}.
  \]
  After tensoring with \(\C((\hbar))\), no nonzero vector in
  \(\mc P((\hbar))\) is killed by a high power of either linear
  Hamiltonian: for a finite nonzero linear combination, every power of
  \(X_{z_1}\) shifts each nonzero basis summand to a distinct
  \(\rho_{a,b+N}\), with nonzero rising-factorial coefficient, and the
  same argument applies to \(X_{z_2}\).  Local nilpotence is invariant
  under module isomorphism.
  Thus any Fourier--Rees comparison that preserves the original
  Hamiltonian action is impossible; the actual positive construction
  changes the polarization and is a flat \( \mathcal D_\hbar \)-module
  bridge only for the Fourier-conjugated action stated above.
\end{proof}

\begin{cor}[Same-action obstruction criterion]
\label{cor:app-matlis-no-flat-same-action-interpolation}
  Let \(M_\hbar\) be a \(\C[[\hbar]]\)-flat family whose generic fibre
  carries the polynomial descendant action of the linear Hamiltonians
  \(z_1,z_2\), so that \(X_{z_1}\) and \(X_{z_2}\) act locally
  nilpotently after tensoring with \(\C((\hbar))\).  Then
  \(M_\hbar((\hbar))\) is not isomorphic, as a module for the same
  Hamiltonian Lie algebra, to \(\mc P((\hbar))\) with the Matlis
  coadjoint action.  The same statement holds for the Rees-normalized
  operators \(\hbar X_{z_i}\), since \(\hbar\) is invertible on the
  generic fibre.  Equivalently, local nilpotence of the two linear
  Hamiltonians on the generic fibre is a necessary condition for a flat
  same-action Rees bridge, and \(\mc P((\hbar))\) violates that condition
  on every nonzero vector.  Any same-action comparison must therefore
  either change the Hamiltonian action by a Fourier twist, enlarge the
  boundary target to a distributional principal-part sector, or state a
  different deformation problem.
\end{cor}

\begin{proof}
  This is the local-nilpotence obstruction of
  Theorem~\ref{thm:app-matlis-polynomial-descendant-obstruction} after extension
  of scalars to \(\C((\hbar))\).  The two linear Hamiltonians remain
  locally nilpotent on the polynomial descendant generic fibre and
  remain non-locally-nilpotent on every nonzero element of
  \(\mc P((\hbar))\).  Multiplication by the invertible scalar
  \(\hbar\) cannot change local nilpotence.
\end{proof}

\begin{rmk}[What the Fourier--Rees construction proves]
\label{rmk:app-matlis-rees-bridge-status}
  The construction above justifies the shared index set
  \((a,b)\) as Fourier--Rees associated-graded bookkeeping, with the factorial and
  \(\hbar^{a+b}\) normalization displayed explicitly.  It also records
  the lattice convention: the Fourier twist gives
  \(\mathcal D_\hbar/(z_1,z_2)\), whose Cech realization is the Rees
  sublattice \(\bigoplus\hbar^{a+b}\C[[\hbar]]\rho_{a,b}\), not the full
  \(\bigoplus\C[[\hbar]]\rho_{a,b}\) before \(\hbar\) is inverted.  It
  does not weaken the no-polynomial theorem.  The theorem says
  precisely where the apparent comparison breaks: Fourier sends the
  Hamiltonian vector fields to Fourier-transformed operators, while the
  Matlis cotangent module uses the original coadjoint vector fields on
  principal parts.
\end{rmk}
